\documentclass{article}
\usepackage{amsmath,amssymb,amsthm,mathtools,mathrsfs}

\begin{document}
\maketitle
\section{Exact fixed-phase observation and its degree cost}
\label{sec:holonomy-degree-cost}

Fix the original packet $h$, integer tensor degree $k\ge1$, and
$\chi=\chi_{h,k}$; retain $g=2\xi$. This section varies only the
polynomial cutoff $N$. The original source density and its measure are
\[
 w_h(u)=\frac{|(g/h)(1/2+iu)|^2}{2\pi},\qquad
 m(u)=w_h^{*k}(u),\qquad d\nu(u)=m(u)\,du,
 \qquad S(u)=k/2+iu .
 \tag{HD1}
\]
The full gamma/eta estimates in the periodized-source proof give
$\int e^{b|u|}w_h(u)\,du<\infty$ for $0<b<\pi/2$.
For this same $b$ the inequality
$e^{b|u_1+\cdots+u_k|}\le\prod_i e^{b|u_i|}$ and Tonelli's theorem
give
$\int e^{b|u|}m(u)\,du\le(\int e^{b|u|}w_h(u)\,du)^k<\infty$.
The measure has no atoms, and its density is positive almost everywhere:
$g/h$ has isolated real-line zeros after the removable selected roots
are filled by their full derivatives. For $k\ge2$ convolution is
positive at every real argument, as proved by integrating the positive
almost-everywhere factors. In particular every finite original
polynomial Gram is positive.

Let $M_N$ be that Gram in the exact coefficient frame
$1,S,\ldots,S^N$. For a fixed real $u_0$ put $S_0=k/2+iu_0$ and
$a_N=(1,S_0,\ldots,S_0^N)$. Retain the actual values
\[
 \Lambda_N(u_0)=a_NM_N^{-1}a_N^*,\qquad
 c_N=M_N^{-1}a_N^*/\Lambda_N(u_0),\qquad
 P_N^{u_0}(S)=\sum_{j=0}^N(c_N)_j S^j .
 \tag{HD2}
\]
The superscript labels this particular polynomial, rather than changing
the original source space $\mathcal P_N$. The Cauchy--Schwarz inequality
in the $M_N$ metric gives
\[
\begin{gathered}
 \Lambda_N(u_0)=\sup_{P\in\mathcal P_N\setminus\{0\}}
    \frac{|P(S_0)|^2}{\int|P(S(u))|^2d\nu(u)},\\
 P_N^{u_0}(S_0)=1,\qquad
 \int|P_N^{u_0}(S(u))|^2d\nu(u)=\Lambda_N(u_0)^{-1}.
\end{gathered}\tag{HD3}
\]
Equality in the supremum follows by inserting $M_N^{-1}a_N^*$.
This calculation retains every original mass and phase of $M_N$.

We prove $\Lambda_N(u_0)\uparrow\infty$, including at a root of $\chi$.
Monotonicity follows from the nested polynomial spaces. Suppose for
contradiction their suprema had finite bound $K$. Then evaluation at
$S_0$ would be a bounded linear functional on the polynomial subspace
of $L^2(\nu)$, and would extend to its closure. The Hilbert space
representation theorem supplies $f$ in that closure with
$P(S_0)=\int\overline f(u)P(S(u))d\nu(u)$ for every polynomial $P$.
Since $S(u)=k/2+iu$ is an invertible affine coordinate over $\mathbb C$,
this includes all polynomials in $u$ with their full coefficients.
Define the finite complex measure
$d\sigma=\overline f\,d\nu-\delta_{u_0}$.
It has all polynomial moments zero. Cauchy--Schwarz gives
\[
 \int e^{b|u|/2}|f(u)|d\nu(u)
 \le\|f\|_{L^2(\nu)}
       \left(\int e^{b|u|}d\nu(u)\right)^{1/2}<\infty.
 \tag{HD4}
\]
Thus $F(z)=\int e^{izu}d\sigma(u)$ is holomorphic in
$|\operatorname{Im}z|<b/2$. On every strictly smaller strip,
dominated differentiation follows by absorbing each power of $|u|$
into the remaining exponential margin. All derivatives of $F$ at zero
vanish. Its power series and the identity theorem imply $F=0$ on
that connected strip, and in particular on the real axis.

For completeness, this forces the measure itself to vanish. For $t>0$
use the Gaussian $\gamma_t(x)=(4\pi t)^{-1/2}e^{-x^2/(4t)}$.
Its integrable Fourier representation is
$\gamma_t(x)=(2\pi)^{-1}\int e^{-t\zeta^2}e^{-i\zeta x}d\zeta$,
obtained by the Gaussian integral. Fubini then gives
$(\sigma*\gamma_t)(x)=0$ since $F(\zeta)=0$.
For each compactly supported continuous $\varphi$,
$\varphi*\gamma_t\to\varphi$ uniformly: split the convolution
into $|y|<\varepsilon$, controlled by uniform continuity, and its
Gaussian tail, whose mass tends to zero. Integration against the
finite measure $\sigma$ therefore gives $\int\varphi\,d\sigma=0$.
The uniqueness of finite regular measures on continuous test functions
gives $\sigma=0$. But $\overline f\nu$ has zero mass at the singleton
$\{u_0\}$ and $\delta_{u_0}$ has mass one there, a contradiction.
This proves
\[
 \Lambda_N(u_0)\longrightarrow\infty,
 \qquad \mathcal U_k\mathcal V P_N^{u_0}\longrightarrow0
       \text{ in the original source norm},\qquad
 P_N^{u_0}(S_0)=1 .
 \tag{HD5}
\]
In particular evaluation on the polynomial domain is not closable as
an operator from its original Hilbert closure to $\mathbb C$: its
graph has the limiting pair $(0,1)$. This conclusion concerns that
specific evaluation map; all finite polynomial maps (HD2) remain exact.

Now fix $L>0$, $0\le\theta<2\pi$, and an integer $n_0$ with
$u_0=(2\pi n_0+\theta)/L$ and $m(u_0)>0$. Its original sampled mass is
$\omega_0=(2\pi/L)m(u_0)>0$. The exact phase source matrix is
\[
 (M_{N,\theta})_{ij}=\frac{2\pi}{L}
 \sum_{n\in\mathbb Z}m((2\pi n+\theta)/L)
 \overline{S((2\pi n+\theta)/L)^i}
 S((2\pi n+\theta)/L)^j .
 \tag{HD6}
\]
Its $n_0$ summand gives
$M_{N,\theta}\succeq\omega_0 a_N^*a_N$, so
\[
 \lambda_{\max}(M_N^{-1}M_{N,\theta})
 \ge\omega_0\Lambda_N(u_0)\longrightarrow\infty .
 \tag{HD7}
\]
For every actual weighted source exponent $a>0$, let
$\kappa_{a,N}=\lambda_{\max}
(M_N^{-1}(M_{a,+,N}+M_{a,-,N}))$.
The proved correlation comparison
$M_{N,\theta}\preceq
(1+\kappa_{a,N}/(e^{aL}-1))M_N$ does not require that its error be
less than one; it follows from summing the same convergent tails.
Consequently
\[
 \boxed{\kappa_{a,N}\ge
 (e^{aL}-1)\bigl(\omega_0\Lambda_N(u_0)-1\bigr),\qquad
 \kappa_{a,N}\longrightarrow\infty .}
 \tag{HD8}
\]
The finite original matrix expression (HD2) is an explicit lower
certificate for this degree cost. To ensure the correlation error
$\kappa_{a,N}/(e^{aL_N}-1)\le\delta_0<1$ by that estimate, the
literal choice must satisfy
$L_N\ge a^{-1}\log(1+\kappa_{a,N}/\delta_0)$.
This derives a constraint on these actual correlation constants at fixed
$h,k$; it does not assert a growth rate with the distinct tensor
parameter $k$, or exclude another arithmetic estimate.

The source isometry over all holonomies is the exact map that relates
this divergent single-phase observation to the full source. Equations
(HG2) and (HD3) give
\[
 \int\|\mathcal Z_{L,\vartheta}\mathcal U_k
               \mathcal V P_N^{u_0}\|^2d\mu(\vartheta)
 =\Lambda_N(u_0)^{-1}\longrightarrow0,
 \qquad
 \|\mathcal Z_{L,\theta}\mathcal U_k
               \mathcal V P_N^{u_0}\|^2\ge\omega_0 .
 \tag{HD9}
\]
The finite columns are regular and define the displayed phase values;
the full direct integral is an equivalence class modulo null phase
sets. Thus these are simultaneous statements about the very same
columns, with the original positive sampled mass retained. In particular
point evaluation in the phase parameter cannot be extended as a bounded
operator from this full source norm through the direct integral.

There is an exact polynomial-quotient criterion as well. Evaluation at
$S_0$ factors through the original $J:\mathbb C[S]\to E$ if and
only if $\chi(S_0)=0$: necessity follows by evaluating $\chi$,
and sufficiency follows from divisibility of every difference in
$\ker J=(\chi)$. In that case the induced map
$E\to\mathbb C$, $[P]\mapsto P(S_0)$, has kernel
$(S-S_0)/(\chi)$, where the notation means the image of the ideal
$(S-S_0)$ in $E$. If $S_0$ has multiplicity $r$ in $\chi$, write
$\chi=(S-S_0)^r q_0$ with $q_0(S_0)\ne0$. The map on its local
factor $\mathbb C[x]/(x^r)$ is the constant coefficient, with kernel
$(x)/(x^r)$; all other primary factors map to zero. The full
nilpotents and every cofactor remain in this map's domain and kernel.
If $\chi(S_0)\ne0$, evaluation is still the exact source map in
(HD2)--(HD9); its value on $\chi$ explains precisely why it does not
factor through $E$.

For $k\ge2$, every $n_0$ has positive mass. For $k=1$, a sampled
selected root of the complete packet has positive mass because, if
$h=(S-S_0)^r h_0$ and $g$ has its complete order $r$ there, then
$(g/h)(S_0)=g^{(r)}(S_0)/(r!h_0(S_0))\ne0$.
Other zero-mass samples are absent from the receiving discrete Hilbert
space and make no claim in (HD7). For the divergence of $\kappa_{a,N}$
alone, one may choose any real $u_0$ where $m(u_0)>0$ and then its
unique phase modulo $2\pi$ for the fixed $L$; such a point exists by
(HD1). This gives (HD8) without asserting that every one-fold phase
has a positive selected root.

Every displayed source, evaluation, quotient and inverse coefficient
map lifts on an original split label $\lambda$ by
$(\lambda,v)\mapsto(\lambda,fv)$ and fixes $\tau$.
For each retained label the exact inverse images under this lift are
\[
 \widetilde f^{-1}(\{(\lambda,0)\})
   =\{(\lambda,v):v\in\ker f\},\qquad
 \widetilde f^{-1}(\{\tau\})=\{\tau\}.
 \tag{HD10}
\]
Thus the nonzero coefficient vectors in the source kernel remain
represented inputs and map to the receiving represented zero. The
external absence has its displayed singleton inverse image.
No identification of a positive sampled mass, a represented zero, or
external absence is used in the degree-cost calculation.

\end{document}
